% !TeX root = ../教材题库_全部条目_转写工作稿.tex
% Source: XJ-B1-C02-EX-P056-Y1002
% 第2章 / 2.3 一元二次不等式 / 教材 p.56、57 / PDF 第63、64页
% OCR未校对：公式与文字均可能有误；未完成数学TeX转写；禁止直接发给学生
% 原题及相邻原解截图（核对用）：../题目原图/XJ-B1-C02-EX-P056-Y1002-01.webp, ../题目原图/XJ-B1-C02-EX-P056-Y1002-02.webp
% 本文件只在转写完成后作为题干使用；OCR 条目尚不是可靠题干。
\begin{bankitem}{2.3 例题 例题8}{第2章 / 2.3 一元二次不等式 / 教材 p.56、57 / PDF 第63、64页}{XJ-B1-C02-EX-P056-Y1002}{OCR未校对：公式与文字均可能有误；未完成数学TeX转写；禁止直接发给学生}
\begin{ocrdraft}
[教材p.56 / PDF第63页]
\par
0河已知汽车从踩刹车到停车所滑行的距离L（m）与速度\ensuremath{{}^\circ\mathrm C}（km/h）之间有如
\par
下关系式：
\par
2
\par
其中差是比例系数，且差>0，M是汽车质量（t）．若一辆卡车不装货物（司机体重忽
\par
略不计）以36km/h的速度行驶时，从刹车到停车需要走20m.当这辆卡车装着等于
\par
车重的货物行驶时，为保证安全，要在发现前面20m处有障碍物时能在离障碍物
\par
5m以外处停车，则最高速度应低于多少（设司机发现障碍物到踩刹车需经过1s)？
\par
解由关系式L一差\ensuremath{\cdot}M\ensuremath{\cdot}，2，将，一36，L一20代人得
\par
362\ensuremath{\cdot}
\par
又卡车司机从发现障碍物到踩刹车需经过1s，这1内卡车行驶的路程为
\par
1000
\par
50
\par
m.因此，
\par
3600
\par
18
\par
50
\par
．一十差．\ensuremath{\cdot}，2<20一5，
\par
18
\par
整理得
\par
[教材p.57 / PDF第64页]
\par
2
\par
，2十一。一3<0．
\par
方程溪十冖3一0有两个不相等的实数根
\par
18
\par
一一27，一1\&
\par
由图象得不等式一，2十一3<0的解集为
\par
18
\par
在这个实际问题中，>0，所以最高速度应低于18km/h．
\end{ocrdraft}
\end{bankitem}
